\section{Open Reasoning}
\sectionkicker{JAPANESE OPEN WEIGHTS}
\begin{wideflow}
{\Large\bfseries 開かれた推論――Holo4の二枚看板とELYZAの和文鍛錬}\par
\smallskip
{\color{SurveyMuted}汎用と和特化の二本立て。重みの素性と許諾、物差しの分母をほどいて読む。}\par
\end{wideflow}

9月28日、H CompanyはHolo4の二品を出した\cite{holo4}。27Bの高密品と35B-A3Bの混合品は、名前が似ていても別物として扱う。高密品は研究・非商用の許諾、混合品はApache-2.0で、ここを取り違えると後の利用の話が崩れる。OSWorldで85.2\%と80.8\%、OSWorld2で61.7\%と30.9\%と置かれるが、私的な評価集合の話は積み残しで、独立の再現もない。重みの素性を確かめない読みはしない。

10月2日、ELYZAはThinking LLM-jp-4の二品を出した\cite{elyza}。33B品が61.69と和文62.18、混合品が58.46と和文59.61で、表の作りは和文に寄せた手法とそろっており、発表側の計測として分母つきで載せる。重みはSHAで留め、許諾はApache-2.0である。再生の軌跡を見る盤面への言及はあるが、再走まではしていない。Holo4の汎用とELYZAの和特化を比べる表は、条件の違いを注記つきで添える。

\begin{claimboundary}[資料と評価の境界]
二品ずつの区別と許諾の違いは利用の前提であり、取り違えない。公開の物差しは発表側の計測であり、私的な評価集合の話は積み残し、独立の再現はない。再生の軌跡を見る盤面の言及は再走の証しではない。
\end{claimboundary}
